% --- UNIVERSAL PREAMBLE BLOCK ---
\documentclass[11pt, a4paper]{article}

\usepackage[utf8]{inputenc}
\usepackage[T1]{fontenc}
\usepackage[spanish, es-tabla, es-nodecimaldot]{babel}
\usepackage[a4paper, top=2.5cm, bottom=2.5cm, left=2.5cm, right=2.5cm]{geometry}

\usepackage{graphicx}
\usepackage{float}
\usepackage{hyperref}
\usepackage{enumitem}
\usepackage{xcolor}
\usepackage{listings}
\usepackage{fancyhdr}
\usepackage{tabularx}
\usepackage{longtable}
\usepackage{xltabular}
\usepackage{booktabs}
\usepackage{colortbl}
\usepackage{caption}
\usepackage{amssymb}

\setlist[itemize]{label=-}
\setcounter{secnumdepth}{0}

\definecolor{codegreen}{rgb}{0,0.6,0}
\definecolor{codegray}{rgb}{0.5,0.5,0.5}
\definecolor{codepurple}{rgb}{0.58,0,0.82}
\definecolor{backcolour}{rgb}{0.96,0.96,0.96}

\lstdefinestyle{pythonstyle}{
    backgroundcolor=\color{backcolour},
    commentstyle=\color{codegreen},
    keywordstyle=\color{magenta},
    numberstyle=\tiny\color{codegray},
    stringstyle=\color{codepurple},
    basicstyle=\ttfamily\small,
    breakatwhitespace=false,
    breaklines=true,
    captionpos=b,
    keepspaces=true,
    numbers=left,
    numbersep=5pt,
    showspaces=false,
    showstringspaces=false,
    showtabs=false,
    tabsize=4,
    language=Python,
    extendedchars=true,
    literate={á}{{\'a}}1 {é}{{\'e}}1 {í}{{\'i}}1 {ó}{{\'o}}1 {ú}{{\'u}}1 {ñ}{{\~n}}1 {Á}{{\'A}}1 {É}{{\'E}}1 {Í}{{\'I}}1 {Ó}{{\'O}}1 {Ú}{{\'U}}1 {Ñ}{{\~N}}1 {π}{{$\pi$}}1
}
\lstset{style=pythonstyle}

\hypersetup{
    colorlinks=true,
    linkcolor=black,
    filecolor=magenta,
    urlcolor=blue,
    pdftitle={Actividad 1 - POO},
}

% Comando auxiliar para dibujar un cuadro de clase estilo UML
\newcommand{\clasebox}[3]{%
\begin{center}
\begin{tabular}{|l|}
\hline
\textbf{#1} \\
\hline
#2
\hline
#3
\hline
\end{tabular}
\end{center}
}

\begin{document}

% ---------------------------------------------------------
% PORTADA
% ---------------------------------------------------------
\begin{titlepage}
    \centering
    \vspace*{1cm}

    {\huge \textbf{Universidad Nacional de Colombia}\par}
    \vspace{0.5cm}
    {\Large Facultad de Minas\par}
    \vspace{1.5cm}

    {\Large \textbf{Programación Orientada a Objetos 2026-2}\par}
    \vspace{2cm}

    {\huge \textbf{Actividad 1}\par}
    \vspace{0.5cm}
    {\large Seguimiento y solución de ejercicios prácticos de Programación Orientada a Objetos\par}

    \vspace{3cm}

    \textbf{Presentado por:}\par
    \vspace{0.5cm}
    {\Large \textbf{Natalia Patricia Remolina Rodríguez}\par}

    \vspace{2cm}

    \textbf{Docente:}\par
    Walter Hugo Arboleda Mazo\par

    \vspace{2cm}

    Medellín, Colombia\par
    17 de septiembre de 2026
\end{titlepage}

% ---------------------------------------------------------
% EJERCICIO 1.4
% ---------------------------------------------------------
\section{Ejercicio 1.4}

\noindent \textbf{Enlace al repositorio en GitHub:} \\
\url{https://github.com/NatoRemolina/UNAL-Programacion-Orientada-a-Objetos/blob/main/actividad-01/ejercicio4.py}

\subsection{Código fuente}
\begin{lstlisting}
import sys

class Familia:
    """
    Clase encargada de calcular y almacenar las edades de los miembros 
    de la familia en función de la edad de Juan.
    """
    def __init__(self):
        # Inicializamos los atributos de las edades en cero
        self.edad_juan = 0
        self.edad_alberto = 0.0
        self.edad_ana = 0.0
        self.edad_mama = 0.0

    def solicitar_datos(self):
        """
        Solicita la edad de Juan por consola. 
        Incluye un manejo de excepciones para asegurar que la entrada sea numérica.
        """
        try:
            print("Ingrese la edad de Juan: ")
            # Lectura de datos desde la entrada estándar
            input_data = sys.stdin.readline()
            
            # Si la entrada está vacía (útil para pruebas automatizadas), asignamos un valor por defecto
            if not input_data.strip():
                self.edad_juan = 9 
                print(f"Usando valor por defecto: {self.edad_juan}")
            else:
                # Convertimos la entrada a un número entero
                self.edad_juan = int(input_data)
                
        except ValueError:
            print("Por favor, ingrese un número entero válido.")

    def calcular_edades(self):
        """
        Aplica las proporciones matemáticas dadas para calcular 
        las edades de Alberto, Ana y la madre en base a la edad de Juan.
        """
        # Edad de Alberto: 2/3 de la edad de Juan
        self.edad_alberto = (2 / 3) * self.edad_juan
        
        # Edad de Ana: 4/3 de la edad de Juan
        self.edad_ana = (4 / 3) * self.edad_juan
        
        # Edad de la mamá: suma de las edades de sus tres hijos
        self.edad_mama = self.edad_juan + self.edad_alberto + self.edad_ana

    def mostrar_resultados(self):
        """
        Imprime en consola las edades calculadas, formateando 
        los números para no mostrar decimales (.0f).
        """
        print(f"La edad de Juan es: {self.edad_juan}")
        print(f"La edad de Alberto es: {self.edad_alberto:.0f}")
        print(f"La edad de Ana es: {self.edad_ana:.0f}")
        print(f"La edad de la Mamá es: {self.edad_mama:.0f}")

# ==========================================
# BLOQUE PRINCIPAL
# ==========================================
if __name__ == "__main__":
    # Instanciamos la clase Familia y ejecutamos los métodos secuencialmente
    familia = Familia()
    familia.solicitar_datos()
    familia.calcular_edades()
    familia.mostrar_resultados()
\end{lstlisting}

\subsection{Diagrama de clases}
\clasebox{Familia}
{%
- edad\_juan \\
- edad\_alberto \\
- edad\_ana \\
- edad\_mama \\
}
{%
+ solicitar\_datos() \\
+ calcular\_edades() \\
+ mostrar\_resultados() \\
}

\newpage
% ---------------------------------------------------------
% EJERCICIO 1.5
% ---------------------------------------------------------
\section{Ejercicio 1.5}

\noindent \textbf{Enlace al repositorio en GitHub:} \\
\url{https://github.com/NatoRemolina/UNAL-Programacion-Orientada-a-Objetos/blob/main/actividad-01/ejercicio5.py}

\subsection{Código fuente}
\begin{lstlisting}
class PruebaEscritorio:
    """
    Clase diseñada para simular paso a paso la prueba de escritorio 
    del algoritmo propuesto, mostrando el estado de las variables en memoria.
    """
    def __init__(self):
        # Inicializamos las variables que se utilizarán en el algoritmo
        self.suma = 0
        self.x = 0
        self.y = 0

    def ejecutar_algoritmo(self):
        """
        Ejecuta secuencialmente las instrucciones del seudocódigo
        e imprime el valor de las variables para evidenciar el seguimiento.
        """
        # Paso 1: Inicializar la variable acumuladora SUMA en 0
        self.suma = 0
        print(f"Paso 1 -> Inicio: SUMA = {self.suma}")

        # Paso 2: Asignar el valor 20 a la variable X
        self.x = 20
        print(f"Paso 2 -> X asignado: {self.x}")

        # Paso 3: Acumular el valor actual de X en SUMA (0 + 20)
        self.suma = self.suma + self.x
        print(f"Paso 3 -> SUMA = SUMA + X: {self.suma}")

        # Paso 4: Asignar el valor 40 a la variable Y
        self.y = 40
        print(f"Paso 4 -> Y asignado: {self.y}")

        # Paso 5: Actualizar X sumándole el cuadrado de Y (20 + 40^2)
        # Nota: En Python, el operador de potenciación es **
        self.x = self.x + (self.y ** 2)
        print(f"Paso 5 -> X = X + Y^2: {self.x}")

        # Paso 6: Acumular en SUMA el resultado de la división X / Y
        self.suma = self.suma + (self.x / self.y)
        print(f"Paso 6 -> SUMA = SUMA + X/Y: {self.suma}")

        # Salida final exigida por el algoritmo
        print("\n--- Resultado ---")
        print(f"EL VALOR DE LA SUMA ES: {self.suma}")

# ==========================================
# BLOQUE PRINCIPAL
# ==========================================
if __name__ == "__main__":
    # Instanciamos la clase y ejecutamos la simulación del algoritmo
    prueba = PruebaEscritorio()
    prueba.ejecutar_algoritmo()
\end{lstlisting}

\subsection{Diagrama de clases}
\clasebox{PruebaEscritorio}
{%
- suma \\
- x \\
- y \\
}
{%
+ ejecutar\_algoritmo() \\
}

\newpage
% ---------------------------------------------------------
% EJERCICIO 1.12
% ---------------------------------------------------------
\section{Ejercicio 1.12}

\noindent \textbf{Enlace al repositorio en GitHub:} \\
\url{https://github.com/NatoRemolina/UNAL-Programacion-Orientada-a-Objetos/blob/main/actividad-01/ejercicio12.py}

\subsection{Código fuente}
\begin{lstlisting}
class Empleado:
    """
    Clase que modela la información laboral de un empleado y calcula 
    su estructura salarial (bruto, retenciones y neto) aplicando
    conceptos de Programación Orientada a Objetos.
    """
    def __init__(self):
        # Inicialización de atributos base del empleado
        self.horas_trabajadas = 0
        self.valor_hora = 0
        self.porcentaje_retencion = 0.0
        
        # Inicialización de atributos que almacenarán los resultados calculados
        self.salario_bruto = 0.0
        self.retencion = 0.0
        self.salario_neto = 0.0

    def ingresar_datos(self):
        """
        Asigna los valores iniciales y constantes definidos por el 
        enunciado del problema.
        """
        self.horas_trabajadas = 48
        self.valor_hora = 5000
        self.porcentaje_retencion = 12.5  # Representa el 12.5% de retención en la fuente

    def calcular_salario(self):
        """
        Ejecuta las operaciones matemáticas para determinar el salario 
        bruto, el monto de descuento y el salario líquido final.
        """
        # Salario bruto: Total ganado sin aplicar ningún descuento
        self.salario_bruto = self.horas_trabajadas * self.valor_hora
        
        # Retención: Se obtiene multiplicando el bruto por el porcentaje dividido entre 100
        self.retencion = self.salario_bruto * (self.porcentaje_retencion / 100)
        
        # Salario neto: La diferencia entre el devengado (bruto) y lo deducido (retención)
        self.salario_neto = self.salario_bruto - self.retencion

    def mostrar_info(self):
        """
        Imprime el desglose salarial en consola, utilizando f-strings 
        para formatear las salidas financieras con separadores de miles y decimales.
        """
        print("--- Información Salarial del Trabajador ---")
        print(f"Horas trabajadas en la semana: {self.horas_trabajadas}")
        print(f"Valor por hora: ${self.valor_hora}")
        print(f"Porcentaje de retención: {self.porcentaje_retencion}%")
        print("-" * 41)
        
        # El formato :,.2f añade comas para miles y muestra exactamente dos decimales
        print(f"Salario Bruto:           ${self.salario_bruto:,.2f}")
        print(f"Retención en la fuente:  ${self.retencion:,.2f}")
        print(f"Salario Neto a pagar:    ${self.salario_neto:,.2f}")

# ==========================================
# BLOQUE PRINCIPAL
# ==========================================
if __name__ == "__main__":
    # Instanciamos el objeto empleado y ejecutamos el flujo de cálculo salarial
    empleado = Empleado()
    empleado.ingresar_datos()
    empleado.calcular_salario()
    empleado.mostrar_info()
\end{lstlisting}

\subsection{Diagrama de clases}
\clasebox{Empleado}
{%
- horas\_trabajadas \\
- valor\_hora \\
- porcentaje\_retencion \\
- salario\_bruto \\
- retencion \\
- salario\_neto \\
}
{%
+ ingresar\_datos() \\
+ calcular\_salario() \\
+ mostrar\_info() \\
}

\newpage
% ---------------------------------------------------------
% EJERCICIO 1.14
% ---------------------------------------------------------
\section{Ejercicio 1.14}

\noindent \textbf{Enlace al repositorio en GitHub:} \\
\url{https://github.com/NatoRemolina/UNAL-Programacion-Orientada-a-Objetos/blob/main/actividad-01/ejercicio14.py}

\subsection{Código fuente}
\begin{lstlisting}
import sys

class Calculadora:
    """
    Clase diseñada para encapsular operaciones matemáticas básicas,
    permitiendo calcular dinámicamente el cuadrado y el cubo de un número.
    """
    def __init__(self):
        # Inicializamos el atributo que almacenará el número a procesar
        self.numero = 0.0

    def leer_numero(self):
        """
        Captura el número ingresado por el usuario a través de la consola.
        Implementa manejo de excepciones para garantizar que el programa
        no colapse ante entradas que no sean numéricas.
        """
        try:
            print("Ingrese un número para calcular su cuadrado y cubo: ")
            # Se utiliza sys.stdin.readline() para una lectura segura en distintos entornos
            input_data = sys.stdin.readline().strip()
            
            if not input_data:
                # Valor por defecto en caso de ejecuciones automatizadas o entradas vacías
                self.numero = 5.0  
                print(f"Usando valor por defecto: {self.numero}")
            else:
                # Convertimos la entrada a un tipo de dato flotante para admitir decimales
                self.numero = float(input_data)
                
        except ValueError:
            print("Entrada inválida. Por favor ingrese un valor numérico. Usando 0.0 por defecto.")
            self.numero = 0.0

    def calcular_potencias(self):
        """
        Calcula las potencias requeridas utilizando el operador de 
        exponenciación de Python (**) e imprime los resultados.
        """
        # El cuadrado es el número elevado a la potencia de 2
        cuadrado = self.numero ** 2
        
        # El cubo es el número elevado a la potencia de 3
        cubo = self.numero ** 3
        
        # Despliegue de los resultados en consola con un formato limpio
        print(f"\nResultados para el número: {self.numero}")
        print("-" * 35)
        print(f"El cuadrado de {self.numero} es: {cuadrado}")
        print(f"El cubo de {self.numero} es: {cubo}")

# ==========================================
# BLOQUE PRINCIPAL
# ==========================================
if __name__ == "__main__":
    # Instanciamos el objeto calculadora y ejecutamos los métodos de lectura y cálculo
    calc = Calculadora()
    calc.leer_numero()
    calc.calcular_potencias()
\end{lstlisting}

\subsection{Diagrama de clases}
\clasebox{Calculadora}
{%
- numero \\
}
{%
+ leer\_numero() \\
+ calcular\_potencias() \\
}

\newpage
% ---------------------------------------------------------
% EJERCICIO 1.17
% ---------------------------------------------------------
\section{Ejercicio 1.17}

\noindent \textbf{Enlace al repositorio en GitHub:} \\
\url{https://github.com/NatoRemolina/UNAL-Programacion-Orientada-a-Objetos/blob/main/actividad-01/ejercicio17.py}

\subsection{Código fuente}
\begin{lstlisting}
import math
import sys

class Circulo:
    """
    Clase que modela un círculo y permite calcular sus propiedades 
    geométricas principales (área y longitud de la circunferencia) 
    a partir de la medida de su radio.
    """
    def __init__(self):
        # Inicializamos los atributos en 0.0 para trabajar con valores decimales (flotantes)
        self.radio = 0.0
        self.area = 0.0
        self.longitud = 0.0

    def leer_radio(self):
        """
        Captura el radio ingresado por el usuario a través de la consola,
        implementando un manejo de excepciones para evitar que el programa 
        falle si se ingresa texto en lugar de números.
        """
        try:
            print("Ingrese el radio del círculo: ")
            # Se usa sys.stdin.readline() para una lectura segura y compatible en distintos entornos
            input_data = sys.stdin.readline().strip()
            
            if not input_data:
                # Se asigna un valor por defecto si la entrada está vacía (útil en pruebas)
                self.radio = 10.0  
                print(f"Usando valor por defecto: {self.radio}")
            else:
                self.radio = float(input_data)
                
        except ValueError:
            print("Entrada inválida. Por favor ingrese un valor numérico. Usando 0.0 por defecto.")
            self.radio = 0.0

    def calcular_propiedades(self):
        """
        Aplica las fórmulas matemáticas estándar haciendo uso de la constante pi 
        proporcionada por la librería math de Python.
        """
        # Cálculo del área: Pi multiplicado por el radio al cuadrado (π * r^2)
        self.area = math.pi * (self.radio ** 2)
        
        # Cálculo de la longitud (perímetro): 2 multiplicado por Pi por el radio (2 * π * r)
        self.longitud = 2 * math.pi * self.radio

    def mostrar_resultados(self):
        """
        Imprime los resultados en consola. Se utiliza el formato :.4f 
        para redondear la salida a exactamente cuatro decimales.
        """
        print(f"\nResultados para el círculo de radio: {self.radio}")
        print("-" * 45)
        print(f"Área del círculo:              {self.area:.4f}")
        print(f"Longitud de la circunferencia: {self.longitud:.4f}")

# ==========================================
# BLOQUE PRINCIPAL
# ==========================================
if __name__ == "__main__":
    # Instanciamos el objeto Circulo y ejecutamos sus métodos en secuencia
    mi_circulo = Circulo()
    mi_circulo.leer_radio()
    mi_circulo.calcular_propiedades()
    mi_circulo.mostrar_resultados()
\end{lstlisting}

\subsection{Diagrama de clases}
\clasebox{Circulo}
{%
- radio \\
- area \\
- longitud \\
}
{%
+ leer\_radio() \\
+ calcular\_propiedades() \\
+ mostrar\_resultados() \\
}

\end{document}
